L'affirmation est fausse. Il suffit d'un contre-exemple : le quadrilatère
non convexe $ABCD$ avec $A(0\,;0)$, $B(2\,;1)$, $C(4\,;0)$ et $D(2\,;3)$
(une \og pointe de flèche \fg). Le sommet $B$ est à l'intérieur du triangle
$ACD$, et la diagonale $[AC]$ passe sous $B$, à l'extérieur du quadrilatère :
son milieu $(2\,;0)$ n'appartient pas au quadrilatère.

L'affirmation devient vraie si l'on se restreint aux polygones
\emph{convexes} : un polygone convexe contient tout segment joignant deux de
ses points, en particulier ses diagonales. Le contre-exemple montre où se
cachait l'hypothèse implicite. \hfill $\blacksquare$
